\switchcolumn

\textit{Custom cmp pbds:}

\begin{lstlisting}
struct cmp {bool operator() (const int &a,const int
&b){return a <= b;}};
typedef
tree<int,null_type,cmp,rb_tree_tag,tree_order_s
tatistics_node_update> ordered_set;
\end{lstlisting}

\switchcolumn

\switchcolumn

\textbf{\textit{GP Hash Table}}
\begin{lstlisting}
struct chash {
  static uint64_t splitmix64(uint64_t x)
  {
    x += 0x9e3779b97f4a7c15;
    x = (x ^ (x >> 30)) * 0xbf58476d1ce4e5b9;
    x = (x ^ (x >> 27)) * 0x94d049bb133111eb;
    return x ^ (x >> 31);
  }
  int operator()(int x) const
  {
    static const uint64_t FIXED_RANDOM =
    chrono::steady_clock::now().time_since_epoch()
    .count();
    return splitmix64(x + FIXED_RANDOM);
  }
};
gp_hash_table<ll, ll, chash> mp;
\end{lstlisting}

\textbf{DSU}
\begin{lstlisting}
struct DSU{
  vector<int> p, rank, setSize;
  int numSets;
  DSU(int N)
  {
    p.assign(N, 0);
    iota(p.begin(), p.end(), 0);
    rank.assign(N, 0), setSize.assign(N, 1);
    numSets = N;
  }
  int findSet(int i) { return (p[i] == i) ? i : (p[i] =
  findSet(p[i])); }
  bool isSameSet(int i, int j) { return findSet(i) ==
  findSet(j); }
  int numDisjointSets() { return numSets; }
  int sizeOfSet(int i) { return setSize[findSet(i)]; }
  void unionSet(int i, int j)
  {
    if (isSameSet(i, j))
      return;
    int x = findSet(i), y = findSet(j);
    if (rank[x] > rank[y])
      swap(x, y);
    p[x] = y;
    if (rank[x] == rank[y])
      ++rank[y];
    setSize[y] += setSize[x];
    --numSets;
  }
};
\end{lstlisting}

\switchcolumn

\textbf{Persistent Segment Tree}
\begin{lstlisting}
struct PST {
  struct Node { int l, r, val; };
  vector<Node> t;
  int T = 0;
  PST(int max_nodes) { // estimate 20 * n for
  // max_nodes
    t.resize(max_nodes);
  }
  int build(int b, int e) {
    int cur = T++;
    if (b == e) return cur;
    int m = (b + e) / 2;
    t[cur].l = build(b, m);
    t[cur].r = build(m + 1, e);
    return cur;
  }
  int update(int pre, int b, int e, int i, int v) {
    int cur = T++;
    t[cur] = t[pre];
    if (b == e) {
      t[cur].val += v;
      return cur;
    }
    int m = (b + e) / 2;
    if (i <= m) t[cur].l = update(t[pre].l, b, m, i,
    v);
    else t[cur].r = update(t[pre].r, m + 1, e, i, v);
    t[cur].val = t[t[cur].l].val + t[t[cur].r].val;
    return cur;
  }
  int query(int pre, int cur, int b, int e, int k) {
    if (b == e) return b;
    int cnt = t[t[cur].l].val - t[t[pre].l].val;
    int m = (b + e) / 2;
    if (cnt >= k) return query(t[pre].l, t[cur].l, b,
    m, k);
    return query(t[pre].r, t[cur].r, m + 1, e, k - cnt);
  }
};
int main() {
  ios::sync_with_stdio(0); cin.tie(0);
  int n, q;
  cin >> n >> q;
  vector<int> a(n + 1), root(n + 1), V;
  map<int, int> comp;
  for (int i = 1; i <= n; i++) cin >> a[i],
  comp[a[i]];
  int id = 0;
  V.resize(comp.size() + 1);
  for (auto& [x, _] : comp) comp[x] = ++id, V[id]
  = x;
  PST pst(20 * (n + 1)); // Allocate enough space
  root[0] = pst.build(1, id);
  for (int i = 1; i <= n; i++)
  root[i] = pst.update(root[i - 1], 1, id,
  comp[a[i]], 1);
  while (q--) {
    int l, r, k;
    cin >> l >> r >> k;
    cout << V[pst.query(root[l - 1], root[r], 1, id,
    k)] << '\n';
  }
  return 0;
}
\end{lstlisting}
%the code returns k-th number in a range l to r if
%the range were sorted
\begin{lstlisting}
int V[N], root[N], a[N];
int32_t main() {
  map<int, int>mp;
  int n, q;
  cin >> n >> q;
  for(int i = 1; i <= n; i++) cin >> a[i], mp[a[i]];
  int c = 0;
  for(auto x : mp) mp[x.first] = ++c, V[c] = x.first;
  root[0] = t.build(1, n);
  for(int i = 1; i <= n; i++) {
    root[i] = t.upd(root[i - 1], 1, n, mp[a[i]], 1);
  }
  while(q--) {
    int l, r, k;
    cin >> l >> r >> k;
    cout << V[t.query(root[l - 1], root[r], 1, n, k)]
    << '\n';
  }
  return 0;
}
\end{lstlisting}

\switchcolumn

\textbf{Fenwick Tree}
Point update and range query.
\begin{lstlisting}
template<typename T>
struct BIT {
  int n;
  vector<T> bit;
  BIT(int n = 0) { init(n); }
  void init(int n_) { n = n_; bit.assign(n + 1, 0); }
  void add(int i, T v) { for (; i <= n; i += i & -i)
    bit[i] += v; }
  T sum(int i) { T r = 0; for (; i > 0; i -= i & -i) r
    += bit[i]; return r; }
  T sum(int l, int r) { return sum(r) - sum(l - 1); }
};
\end{lstlisting}
Usage: \texttt{BIT<ll> fenw(5);}

\textit{//Fenwick 2D}
\begin{lstlisting}
struct BIT2D {
  int n, m;
  vector<vector<ll>> bit;
  BIT2D(int n, int m): n(n), m(m) {
    bit.assign(n + 1, vector<ll>(m + 1, 0));
  }
  void add(int x, int y, ll val) {
    for (int i = x; i <= n; i += i & -i)
      for (int j = y; j <= m; j += j & -j)
        bit[i][j] += val;
  }
  ll sum(int x, int y) {
    ll res = 0;
    for (int i = x; i > 0; i -= i & -i)
      for (int j = y; j > 0; j -= j & -j)
        res += bit[i][j];
    return res;
  }
  ll query(int x1, int y1, int x2, int y2) {
    return sum(x2, y2) - sum(x1 - 1, y2) - sum(x2,
    y1 - 1) + sum(x1 - 1, y1 - 1);
  }
};
\end{lstlisting}
\begin{lstlisting}
BIT2D fenw2(4, 4);
fenw2.add(1, 1, 5);
fenw2.add(3, 2, 7);
cout << fenw2.query(1, 1, 3, 3) << "\n"; // sum
// inside rectangle [(1,1)-(3,3)]
\end{lstlisting}

Range update and range query.
\begin{lstlisting}
using ll = long long;
ll bit1[200005] ,bit2[200005],n,m;
void update(ll bit[],ll i,ll val){
  for(;i<=n;i+=(i&(-i))){
    bit[i]+=val;
  }
}
void range_update(ll l, ll r,ll val){
  update(bit1,l,val);
  update(bit1,r+1,-val);
  update(bit2,l,val*(l-1));
  update(bit2,r+1,-(r*val));
}
ll sum(ll bit[],ll i){
  ll re = 0;
  for(;i>0;i-=(i&(-i))){
    re+=bit[i];
  }
  return re;
}
ll prefix_sum(ll i){
  return sum(bit1,i)*i-sum(bit2,i);
}
ll range_sum(ll l,ll r){
  return prefix_sum(r)-prefix_sum(l-1);
}
\end{lstlisting}
\texttt{range\_update(l,r,val);}\\
\texttt{range\_sum(l,r);}

\switchcolumn

\textbf{Sparse Table}
Max/Min/GCD Query
\begin{lstlisting}
struct SparseTable {
  SparseTable(vector<int>& data)
  {
    int n = data.size();
    int maxLog = log2(n) + 1;
    table.assign(n, vector<int>(maxLog));
    for (int i = 0; i < n; ++i) {
      table[i][0] = data[i];
    }
    for (int j = 1; (1 << j) <= n; ++j) {
      for (int i = 0; (i + (1 << j) - 1) < n; ++i) {
        table[i][j] = max(table[i][j - 1], table[i +
        (1 << (j - 1))][j - 1]);
      }
    }
  }
  int query(int L, int R)
  {
    int j = log2(R - L + 1);
    return max(table[L][j], table[R - (1 << j) +
    1][j]);
  }
  vector<vector<int>> table;
};
\end{lstlisting}
\begin{lstlisting}
int main()
{
  vector<int> data = { 1, 3, 2, 7, 9, 11, 3, 5, 6, 8 };
  SparseTable st(data);
  cout << "Max in range [1, 4]: " << st.query(1, 4)
  << endl;
}
\end{lstlisting}

\textbf{Sparse Table 2D} (up to 500 x 500 is safe)
\begin{lstlisting}
struct SparseTable2D {
  static const int LG = 10;
  int n, m;
  vector<vector<int>> a;
  vector<int> lg2;
  vector<vector<vector<vector<int>>>> st;
  SparseTable2D(int n, int m) : n(n), m(m) {
    a.assign(n, vector<int>(m, 0));
    lg2.resize(max(n, m) + 1, 0);
    for (int i = 2; i <= max(n, m); i++) lg2[i] =
    lg2[i >> 1] + 1;
    st.assign(n, vector<vector<vector<int>>>(m,
    vector<vector<int>>(LG, vector<int>(LG,
    0))));
  }
  void set_val(int i, int j, int val) { a[i][j] = val; }
  void build() {
    for (int i = 0; i < n; i++)
      for (int j = 0; j < m; j++)
        st[i][j][0][0] = a[i][j];
    for (int a = 0; a < LG; a++) {
      for (int b = 0; b < LG; b++) {
        if (a + b == 0) continue;
        for (int i = 0; i + (1 << a) <= n; i++) {
          for (int j = 0; j + (1 << b) <= m; j++) {
            if (!a)
              st[i][j][a][b] = max(st[i][j][a][b-1],
              st[i][j + (1 << (b - 1))][a][b - 1]);
            else
              st[i][j][a][b] = max(st[i][j][a-1][b],
              st[i + (1 << (a - 1))][j][a - 1][b]);
          }
        }
      }
    }
  }
  int query(int x1, int y1, int x2, int y2) {
    int a = lg2[x2 - x1 + 1], b = lg2[y2 - y1 + 1];
    int m1 = max(st[x1][y1][a][b], st[x2 - (1 << a)
    + 1][y1][a][b]);
    int m2 = max(st[x1][y2 - (1 << b) + 1][a][b],
    st[x2 - (1 << a) + 1][y2 - (1 << b) + 1][a][b]);
    return max(m1, m2);
  }
};
\end{lstlisting}
\begin{lstlisting}
// st.set_val(i, j, grid[i][j]);
// st.build(); // st.query(0, 0, 1, 2); // max in
// submatrix [(0,0) to (1,2)]
\end{lstlisting}

\switchcolumn

\textbf{Merge Sort Tree}
\begin{lstlisting}
struct MergeSortTree {
  vector<vector<int>> tree;
  vector<int> arr;
  int n;
  MergeSortTree(vector<int>& a)
  {
    arr = a, n = a.size();
    tree.resize(4 * n);
    build(1, 0, n - 1);
  }
  void build(int node, int l, int r)
  {
    if (l == r) {
      tree[node].push_back(arr[l]);
      return;
    }
    int mid = (l + r) / 2;
    build(node * 2, l, mid), build(node * 2 + 1,
    mid + 1, r);
    merge(tree[node * 2].begin(), tree[node *
    2].end(), tree[node * 2 + 1].begin(), tree[node * 2
    + 1].end(), back_inserter(tree[node]));
  }
  int query(int node, int l, int r, int ql, int qr, int
  k)
  {
    if (l > qr || r < ql)
      return 0;
    if (l >= ql && r <= qr) {
      int hi = tree[node].size() - 1, lo = 0, ind =
      -1;
      // while (hi >= lo) { } // Modify your
      // function here
      return ind;
    }
    int mid = (l + r) / 2;
    return query(node * 2, l, mid, ql, qr, k) +
    query(node * 2 + 1, mid + 1, r, ql, qr, k);
  }
};
\end{lstlisting}

Find the number of distinct elements in a range.
\begin{lstlisting}
void merge(vector<int> tree[], int treeNode)
{
  int len1 = tree[2 * treeNode].size();
  int len2 = tree[2 * treeNode + 1].size();
  int index1 = 0, index2 = 0;
  while (index1 < len1 && index2 < len2) {
    if (tree[2 * treeNode][index1]
    > tree[2 * treeNode + 1][index2]) {
      tree[treeNode].push_back(
      tree[2 * treeNode + 1][index2]);
      index2++;
    }
    else {
      tree[treeNode].push_back(
      tree[2 * treeNode][index1]);
      index1++;
    }
  }
  while (index1 < len1) {
    tree[treeNode].push_back(
    tree[2 * treeNode][index1]);
    index1++;
  }
  while (index2 < len2) {
    tree[treeNode].push_back(
    tree[2 * treeNode + 1][index2]);
    index2++;
  }
  return;
}
\end{lstlisting}

\begin{lstlisting}
void build(vector<int> tree[], int* arr, int start,
int end,
int treeNode)
{
  if (start == end) {
    tree[treeNode].push_back(arr[start]);
    return;
  }
  int mid = (start + end) / 2;
  build(tree, arr, start, mid, 2 * treeNode);
  build(tree, arr, mid + 1, end, 2 * treeNode + 1);
  merge(tree, treeNode);
  return;
}
\end{lstlisting}

\switchcolumn

\begin{lstlisting}
int query(vector<int> tree[], int treeNode, int
start,
int end, int left, int right)
{
  if (start > right || end < left) {
    return 0;
  }
  if (start >= left && end <= right) {
    return tree[treeNode].end()
    - upper_bound(tree[treeNode].begin(),
    tree[treeNode].end(), right);
  }
  int mid = (start + end) / 2;
  int op1 = query(tree, 2 * treeNode, start, mid,
  left,
  right);
  int op2 = query(tree, 2 * treeNode + 1, mid + 1,
  end,
  left, right);
  return op1 + op2;
}
int main()
{
  int n; cin >> n;
  int arr[n];
  for(int i=0;i<n;i++) cin >> arr[i];
  int next_right[n];
  vector<int> tree[4 * n];
  unordered_map<int, int> ump;
  for (int i = n - 1; i >= 0; i--) {
    if (ump[arr[i]] == 0) {
      next_right[i] = n;
      ump[arr[i]] = i;
    }
    else {
      next_right[i] = ump[arr[i]];
      ump[arr[i]] = i;
    }
  }
  build(tree, next_right, 0, n - 1, 1);
  int q; cin >> q;
  while(q--){
    int l, r; cin >> l >> r;
    l-- , r--;
    cout << query(tree,1,0,n-1,l,r) << "\n";
  }
  return 0;
}
\end{lstlisting}

\textbf{Kth smallest value in a range (SegTree)}
\begin{lstlisting}
#define N (int)1e5
vector<int> seg[N];
void build(int ci, int st, int end,
pair<int, int>* B)
{
  if (st == end) {
    seg[ci].push_back(B[st].second);
    return;
  }
  int mid = (st + end) / 2;
  build(2 * ci + 1, st, mid, B);
  build(2 * ci + 2, mid + 1, end, B);
  merge(seg[2 * ci + 1].begin(), seg[2 * ci +
  1].end(),
  seg[2 * ci + 2].begin(), seg[2 * ci + 2].end(),
  back_inserter(seg[ci]));
}
\end{lstlisting}

\switchcolumn

\begin{lstlisting}
int query(int ci, int st, int end,
int l, int r, int k)
{
  if (st == end)
    return seg[ci][0];
  int p = upper_bound(seg[2 * ci + 1].begin(),
  seg[2 * ci + 1].end(), r)
  - lower_bound(seg[2 * ci + 1].begin(),
  seg[2 * ci + 1].end(), l);
  int mid = (st + end) / 2;
  if (p >= k)
    return query(2 * ci + 1, st, mid, l, r, k);
  else
    return query(2 * ci + 2, mid + 1, end, l, r, k-p);
}
int main()
{
  int arr[] = { 3, 1, 5, 2, 4, 7, 8, 6 };
  int n = sizeof(arr) / sizeof(arr[0]);
  pair<int, int> B[n];
  for (int i = 0; i < n; i++) {
    B[i] = { arr[i], i };
  }
  sort(B, B + n);
  build(0, 0, n - 1, B);
  cout << "3rd smallest element in range 3 to 7 is:
  "
  << arr[query(0, 0, n - 1, 2, 6, 3)] << "\n";
}
